\section{Configuration reference}
\label{sec:appendix-configuration}

A PostgreSQL configuration setting is also called a \emph{GUC}. This reference
covers all 75 production BioSeqDB settings: 73 in the main extension and two in
\texttt{bioseq\_cache}. Three additional settings belong to test builds and appear
separately at the end. These are the defaults registered by the implementation
studied in this thesis. Server, database, role and session settings can replace
them. A backend is the PostgreSQL server process executing the work.

\subsection{Reading and changing a setting}

\begin{itemize}
  \item \texttt{SHOW} reads one setting. The \texttt{pg\_settings} view also
    reports its current value, type, range, allowed enum values and change
    context. Extension settings become available when the extension library is
    loaded. A missing row is therefore not proof that a setting is unsupported.
  \item \texttt{SET} changes the current session. \texttt{RESET} restores that
    session's reset value, which can differ from the built-in default below.
  \item \texttt{SET LOCAL} applies inside the current transaction and ends at
    commit or rollback. It is useful when comparing search choices without
    leaving a changed session behind. A rolled-back \texttt{SET} is also undone.
\end{itemize}
These commands follow PostgreSQL's configuration rules
\autocite{postgresql17_set,postgresql17_pg_settings}.

\par\noindent\begin{minipage}{\linewidth}
\begin{lstlisting}[style=bioseq-appendix-sql,
  caption={Inspect a default and limit a change to one transaction.},
  label={lst:appendix-guc-session}]
SHOW bioseq.aligner.minimizer.algorithm;
SET bioseq.aligner.minimizer.algorithm = 'scalar';
RESET bioseq.aligner.minimizer.algorithm;

BEGIN;
SET LOCAL bioseq.aligner.minimizer.algorithm = 'scalar';
-- Run the search to be measured here.
COMMIT;

SELECT name, setting, unit, boot_val, reset_val,
       context, min_val, max_val, enumvals
FROM pg_settings
WHERE name LIKE 'bioseq.%'
   OR name LIKE 'bioseq_cache.%'
ORDER BY name;
\end{lstlisting}
\end{minipage}

A setting marked \textbf{hidden} is omitted from both \texttt{SHOW ALL} and
\texttt{pg\_settings}. It can still be read by its exact name with \texttt{SHOW}
\autocite{postgresql17_guc_functions}.
Hidden does not mean unimplemented. Conversely, accepting a name before the
extension loads does not prove that it is a registered BioSeqDB option.

\subsection{How to read the tables}

Each table states a prefix. Join that prefix, a dot, and the row's suffix to
obtain the complete name. For example, prefix \texttt{bioseq.column\_cache}
and suffix \texttt{enabled} mean \texttt{bioseq.column\_cache.enabled}.
\begin{itemize}
  \item The first column gives the suffix and default. All settings can be
    changed by the current user unless marked \textbf{superuser} or
    \textbf{reload}. A superuser setting also permits a role granted the
    corresponding PostgreSQL \texttt{SET} privilege. Reload settings are
    changed through server configuration and a configuration reload, not
    ordinary session \texttt{SET}.
  \item Numeric intervals include both endpoints. \emph{Integer} settings store whole
    numbers, while \emph{real} settings also store fractions. Boolean settings accept
    \texttt{on} and \texttt{off}. Enum spellings are listed explicitly.
  \item Values in MiB use PostgreSQL's \texttt{MB} unit, equal to $2^{20}$ bytes.
    The shard setting takes a plain integer GiB count, equal to $2^{30}$ bytes.
    A unit described in the text is not necessarily a registered GUC unit:
    counts of symbols, anchors, rows and bits take plain numbers.
  \item The ranges describe what \texttt{SET} accepts. They do not guarantee
    that every combination works with every algorithm. The notes below each
    group state the relevant restrictions.
\end{itemize}
See PostgreSQL's parameter documentation for units and change contexts
\autocite{postgresql17_setting_parameters,postgresql17_pg_settings}.

The search options object can replace corresponding defaults for one
\texttt{mmsearch} call. A supplied non-NULL field replaces that default.
Section~\ref{sec:appendix-search-options} lists the mapping, the fields without
GUC equivalents, and when a scan reads the defaults. Algorithm-affecting GUC
changes reset cached plans where required by the extension.

\Needspace{16\baselineskip}
\subsection{Physical storage defaults}

Prefix: \nolinkurl{bioseq.storage}.

\begingroup
\small
\setlength{\tabcolsep}{4pt}
\renewcommand{\arraystretch}{1.16}
\begin{longtable}{@{}>{\raggedright\arraybackslash}p{0.34\linewidth}>{\raggedright\arraybackslash}p{0.22\linewidth}>{\raggedright\arraybackslash}p{0.40\linewidth}@{}}
\toprule
Setting and default & Accepted values & Effect \\
\midrule
\endfirsthead
\multicolumn{3}{@{}l}{Prefix: \nolinkurl{bioseq.storage} (continued)} \\
\toprule
Setting and default & Accepted values & Effect \\
\midrule
\endhead
\bottomrule
\endfoot
% GUC: bioseq.storage.compression.algorithm
\nolinkurl{compression.algorithm}\par Default: best available & Enum\par \nolinkurl{none},\allowbreak\ \nolinkurl{pglz},\allowbreak\ \nolinkurl{lz4},\allowbreak\ \nolinkurl{zstd} & Codec for compressed column chunks. LZ4 and Zstandard require compiled support. \\[4pt]
% GUC: bioseq.storage.compression.level
\nolinkurl{compression.level}\par Default: 3 & Integer\par $1\ldots 19$ & Sets the compression level used for compressed column chunks. \\[4pt]
% GUC: bioseq.storage.stripe.row_limit
\nolinkurl{stripe.row_limit}\par Default: 150000 & Integer\par $1000\ldots 100000000$ & Limits the number of rows stored in one stripe. \\[4pt]
% GUC: bioseq.storage.chunk_group.row_limit
\nolinkurl{chunk_group.row_limit}\par Default: 10000 & Integer\par $1000\ldots 100000000$ & Limits the number of rows stored in one chunk group. \\[4pt]
% GUC: bioseq.storage.sequence_shard.max_size_gib
\nolinkurl{sequence_shard.max_size_gib}\par Default: 1024\par \textbf{superuser} & Integer\par $0\ldots 2147483647$ & Maximum sequence-shard size in GiB. Zero starts a new shard for each inserted sequence. \\[4pt]
% GUC: bioseq.storage.dml.enabled
\nolinkurl{dml.enabled}\par Default: \texttt{on}\par \textbf{hidden} & Boolean\par \texttt{off}, \texttt{on} & Applies deletion masks during reads. Keep on for ordinary use. \\[4pt]
\end{longtable}
\endgroup

Compression and row limits become stored table options when a table is created.
Changing a default does not rewrite existing chunks or replace those stored
options. The compression default is \texttt{zstd} when available, otherwise
\texttt{lz4}, otherwise \texttt{pglz}. \texttt{none} disables compression.
The compression level affects Zstandard only in the current implementation.

A sequence must fit wholly inside one positive-sized shard. Changing the shard
limit does not move existing sequences. Keep \texttt{dml.enabled} on: it controls
deleted-row filtering and is not a read-only mode.

\Needspace{16\baselineskip}
\subsection{Decompressed column cache}

Prefix: \nolinkurl{bioseq.column_cache}.

\begingroup
\small
\setlength{\tabcolsep}{4pt}
\renewcommand{\arraystretch}{1.16}
\begin{longtable}{@{}>{\raggedright\arraybackslash}p{0.34\linewidth}>{\raggedright\arraybackslash}p{0.22\linewidth}>{\raggedright\arraybackslash}p{0.40\linewidth}@{}}
\toprule
Setting and default & Accepted values & Effect \\
\midrule
\endfirsthead
\multicolumn{3}{@{}l}{Prefix: \nolinkurl{bioseq.column_cache} (continued)} \\
\toprule
Setting and default & Accepted values & Effect \\
\midrule
\endhead
\bottomrule
\endfoot
% GUC: bioseq.column_cache.enabled
\nolinkurl{enabled}\par Default: \texttt{off} & Boolean\par \texttt{off}, \texttt{on} & Caches decompressed column chunks within a backend. \\[4pt]
% GUC: bioseq.column_cache.capacity_mb
\nolinkurl{capacity_mb}\par Default: 200 & Integer\par $20\ldots 20000$ & Capacity of the backend-local decompressed column-chunk cache. Unit: MiB. \\[4pt]
\end{longtable}
\endgroup

This backend-local cache holds decompressed buffers from compressed columns.
It differs from PostgreSQL shared buffers and the query-result cache. Buffers
currently in use are protected from eviction. Normal scan teardown clears this
cache, so the capacity is neither a promise of reuse across statements nor a
limit on total backend memory.

\Needspace{16\baselineskip}
\subsection{Planning and SQL execution}

Prefix: \nolinkurl{bioseq}.

\begingroup
\small
\setlength{\tabcolsep}{4pt}
\renewcommand{\arraystretch}{1.16}
\begin{longtable}{@{}>{\raggedright\arraybackslash}p{0.34\linewidth}>{\raggedright\arraybackslash}p{0.22\linewidth}>{\raggedright\arraybackslash}p{0.40\linewidth}@{}}
\toprule
Setting and default & Accepted values & Effect \\
\midrule
\endfirsthead
\multicolumn{3}{@{}l}{Prefix: \nolinkurl{bioseq} (continued)} \\
\toprule
Setting and default & Accepted values & Effect \\
\midrule
\endhead
\bottomrule
\endfoot
% GUC: bioseq.vectorization.enabled
\nolinkurl{vectorization.enabled}\par Default: \texttt{on}\par \textbf{hidden} & Boolean\par \texttt{off}, \texttt{on} & Enables supported SQL filter and aggregate vectorization. \\[4pt]
% GUC: bioseq.execution.parallel.enabled
\nolinkurl{execution.parallel.enabled}\par Default: \texttt{on}\par \textbf{hidden} & Boolean\par \texttt{off}, \texttt{on} & Allows execution to use parallel workers. \\[4pt]
% GUC: bioseq.execution.parallel.minimum_processes
\nolinkurl{execution.parallel.minimum_processes}\par Default: 8 & Integer\par $1\ldots 32$ & Requested PostgreSQL worker count for a BioSeqDB parallel scan path. \\[4pt]
% GUC: bioseq.planner.index_scan.enabled
\nolinkurl{planner.index_scan.enabled}\par Default: \texttt{on}\par \textbf{hidden} & Boolean\par \texttt{off}, \texttt{on} & Enables Bioseq index scan paths. \\[4pt]
% GUC: bioseq.planner.custom_scan.enabled
\nolinkurl{planner.custom_scan.enabled}\par Default: \texttt{on}\par \textbf{hidden} & Boolean\par \texttt{off}, \texttt{on} & Enables custom scan paths for Bioseq tables. \\[4pt]
% GUC: bioseq.planner.custom_scan.max_paths
\nolinkurl{planner.custom_scan.max_paths}\par Default: 64\par \textbf{hidden} & Integer\par $1\ldots 1024$ & Bounds the parameterized predicate alternatives explored for custom-scan paths. \\[4pt]
% GUC: bioseq.planner.qual_pushdown.enabled
\nolinkurl{planner.qual_pushdown.enabled}\par Default: \texttt{on}\par \textbf{hidden} & Boolean\par \texttt{off}, \texttt{on} & Pushes supported filter predicates into custom scans. \\[4pt]
% GUC: bioseq.planner.qual_pushdown.correlation_threshold
\nolinkurl{planner.qual_pushdown.correlation_threshold}\par Default: 0.4\par \textbf{hidden} & Real\par $0\ldots 1$ & Minimum absolute column correlation used by the predicate-pushdown decision. \\[4pt]
% GUC: bioseq.planner.diagnostics.log_level
\nolinkurl{planner.diagnostics.log_level}\par Default: \nolinkurl{debug3} & Enum\par \nolinkurl{debug5},\allowbreak\ \nolinkurl{debug4},\allowbreak\ \nolinkurl{debug3},\allowbreak\ \nolinkurl{debug2},\allowbreak\ \nolinkurl{debug1},\allowbreak\ \nolinkurl{log},\allowbreak\ \nolinkurl{info},\allowbreak\ \nolinkurl{notice},\allowbreak\ \nolinkurl{warning} & Sets the message level for planner diagnostics. \\[4pt]
% GUC: bioseq.planner.sequence_sort_scan.enabled
\nolinkurl{planner.sequence_sort_scan.enabled}\par Default: \texttt{off}\par \textbf{hidden} & Boolean\par \texttt{off}, \texttt{on} & Enables the optional sequence batch-sort scan for eligible single sequence sort keys. \\[4pt]
% GUC: bioseq.diagnostics.logging.enabled
\nolinkurl{diagnostics.logging.enabled}\par Default: \texttt{off}\par \textbf{superuser} & Boolean\par \texttt{off}, \texttt{on} & Enables extension diagnostic logging. \\[4pt]
\end{longtable}
\endgroup

SQL vectorization applies to supported filters and aggregates. It is separate
from the SIMD minimizer and chaining choices. Lazy sequence projections can
keep scalar SQL execution even when vectorization is enabled.

The parallel-process setting requests a planner worker count. It does not
reserve that many PostgreSQL processes. Actual worker launch depends on
PostgreSQL. The sequence-sort scan is an optional prototype for eligible single
sequence sort keys.

Do not disable custom scans in a normal \texttt{mmsearch} query. The search
predicate depends on the custom-scan rewrite and is not an ordinary function
that performs a slower fallback search.

\Needspace{16\baselineskip}
\subsection{Task scheduling}

Prefix: \nolinkurl{bioseq.scheduler}.

\begingroup
\small
\setlength{\tabcolsep}{4pt}
\renewcommand{\arraystretch}{1.16}
\begin{longtable}{@{}>{\raggedright\arraybackslash}p{0.34\linewidth}>{\raggedright\arraybackslash}p{0.22\linewidth}>{\raggedright\arraybackslash}p{0.40\linewidth}@{}}
\toprule
Setting and default & Accepted values & Effect \\
\midrule
\endfirsthead
\multicolumn{3}{@{}l}{Prefix: \nolinkurl{bioseq.scheduler} (continued)} \\
\toprule
Setting and default & Accepted values & Effect \\
\midrule
\endhead
\bottomrule
\endfoot
% GUC: bioseq.scheduler.backend.max_concurrency
\nolinkurl{backend.max_concurrency}\par Default: 0\par \textbf{reload} & Integer\par $0\ldots 1024$ & Maximum participating task slots in each backend. Zero uses the detected CPU count. \\[4pt]
% GUC: bioseq.scheduler.pipeline.mode
\nolinkurl{pipeline.mode}\par Default: \nolinkurl{single_thread} & Enum\par \nolinkurl{single_thread},\allowbreak\ \nolinkurl{parallel_dispatch},\allowbreak\ \nolinkurl{task_parallel} & Chooses caller-only execution, eligible parallel dispatch, or parallel work within the pipeline. \\[4pt]
% GUC: bioseq.scheduler.postgres_parallel.enabled
\nolinkurl{postgres_parallel.enabled}\par Default: \texttt{off} & Boolean\par \texttt{off}, \texttt{on} & Allows a PostgreSQL parallel-worker process to use more than one oneTBB task slot. \\[4pt]
\end{longtable}
\endgroup

\begin{itemize}
  \item \texttt{single\_thread} executes dispatch and coarse pipeline work on
    the caller.
  \item \texttt{parallel\_dispatch} permits independent eligible dispatch jobs
    to run together. Internal diagnostic output can spell this mode
    \texttt{dispatch\_parallel}.
  \item \texttt{task\_parallel} also permits parallel regions inside the
    pipeline. Small batches, profiling and sequence storage that workers cannot
    access safely can still require serial work.
\end{itemize}
The concurrency cap applies to each backend and includes a participating
caller. It is not a cluster quota or a count of extra PostgreSQL workers.
PostgreSQL parallel-worker processes use one task slot unless the last setting
is enabled.

\Needspace{16\baselineskip}
\subsection{Persistent index maintenance and reads}

Prefix: \nolinkurl{bioseq.index}.

\begingroup
\small
\setlength{\tabcolsep}{4pt}
\renewcommand{\arraystretch}{1.16}
\begin{longtable}{@{}>{\raggedright\arraybackslash}p{0.34\linewidth}>{\raggedright\arraybackslash}p{0.22\linewidth}>{\raggedright\arraybackslash}p{0.40\linewidth}@{}}
\toprule
Setting and default & Accepted values & Effect \\
\midrule
\endfirsthead
\multicolumn{3}{@{}l}{Prefix: \nolinkurl{bioseq.index} (continued)} \\
\toprule
Setting and default & Accepted values & Effect \\
\midrule
\endhead
\bottomrule
\endfoot
% GUC: bioseq.index.insert_maintenance.enabled
\nolinkurl{insert_maintenance.enabled}\par Default: \texttt{on} & Boolean\par \texttt{off}, \texttt{on} & Allows insert-time index maintenance according to each index's stored update policy. \\[4pt]
% GUC: bioseq.index.sequence_read.enabled
\nolinkurl{sequence_read.enabled}\par Default: \texttt{on} & Boolean\par \texttt{off}, \texttt{on} & Allows compatible persisted per-sequence indexes to be read. \\[4pt]
% GUC: bioseq.index.async_persistence.enabled
\nolinkurl{async_persistence.enabled}\par Default: \texttt{on} & Boolean\par \texttt{off}, \texttt{on} & Allows eligible newly built minimizer indexes to be staged for later persistence. \\[4pt]
% GUC: bioseq.index.recompute_async_persistence.enabled
\nolinkurl{recompute_async_persistence.enabled}\par Default: \texttt{off} & Boolean\par \texttt{off}, \texttt{on} & Allows eligible recomputed minimizer indexes to use asynchronous persistence. \\[4pt]
% GUC: bioseq.index.memory_cap_mb
\nolinkurl{memory_cap_mb}\par Default: 0 & Integer\par $0\ldots 2147483647$ & Cap for accounted resident-index memory reservations. Zero uses the runtime memory policy. Unit: MiB. \\[4pt]
% GUC: bioseq.index.sequence_minimizer_sort.algorithm
\nolinkurl{sequence_minimizer_sort.algorithm}\par Default: \nolinkurl{std_sort}\par \textbf{hidden} & Enum\par \nolinkurl{std_sort},\allowbreak\ \nolinkurl{simd_u128_vqsort},\allowbreak\ \nolinkurl{simd_u128_vec_stream_1},\allowbreak\ \nolinkurl{simd_u128_vec_stream_2},\allowbreak\ \nolinkurl{simd_u128_vec_stream_4} & Selects the sorter used for minimizer sequence pairs. \\[4pt]
\end{longtable}
\endgroup

\begin{itemize}
  \item \texttt{insert\_maintenance.enabled} is a session gate. Each index
    separately stores its \texttt{UPDATE\_MANUALLY} or
    \texttt{UPDATE\_ON\_INSERT} policy.
  \item Turning \texttt{sequence\_read.enabled} off prevents stored-index
    reuse. It does not disable index maintenance or later writes.
  \item Recompute persistence requires both asynchronous-persistence switches
    on, a real source relation, an eligible index policy and a supported
    minimizer scheme. This path currently requires $k=15,w=10$ and excludes
    bounded ntHash.
  \item Async staging is best effort. Jobs become available to the persistence
    worker after commit. Abort releases them. Missing queue resources can
    prevent persistence, so a successful search does not promise a newly
    stored index.
\end{itemize}
The index layer counts its memory reservations against this cap and the
available memory budget, using whichever limit is smaller. Other query buffers
are outside this count, so the setting does not limit total query memory.

\Needspace{16\baselineskip}
\subsection{Index construction and loading}

Prefix: \nolinkurl{bioseq}.

\begingroup
\small
\setlength{\tabcolsep}{4pt}
\renewcommand{\arraystretch}{1.16}
\begin{longtable}{@{}>{\raggedright\arraybackslash}p{0.34\linewidth}>{\raggedright\arraybackslash}p{0.22\linewidth}>{\raggedright\arraybackslash}p{0.40\linewidth}@{}}
\toprule
Setting and default & Accepted values & Effect \\
\midrule
\endfirsthead
\multicolumn{3}{@{}l}{Prefix: \nolinkurl{bioseq} (continued)} \\
\toprule
Setting and default & Accepted values & Effect \\
\midrule
\endhead
\bottomrule
\endfoot
% GUC: bioseq.index_build.run_buffer_bytes
\nolinkurl{index_build.run_buffer_bytes}\par Default: 0\par \textbf{superuser, hidden} & Integer\par $0\ldots 2147483647$ & Overrides the external-sort run buffer size. Zero uses the computed production size. Unit: bytes. \\[4pt]
% GUC: bioseq.load.write_overlap.enabled
\nolinkurl{load.write_overlap.enabled}\par Default: \texttt{off} & Boolean\par \texttt{off}, \texttt{on} & Overlaps input decoding with column writes during loading. \\[4pt]
\end{longtable}
\endgroup

The run-buffer override is an administrative tool for external sorting.
Zero leaves its size to the computed production policy. Loading overlap places
input decoding and PostgreSQL writes on opposite sides of a bounded queue.
The PostgreSQL backend still performs the table writes.

\Needspace{16\baselineskip}
\subsection{Minimizer extraction}

Prefix: \nolinkurl{bioseq.aligner.minimizer}.

\begingroup
\small
\setlength{\tabcolsep}{4pt}
\renewcommand{\arraystretch}{1.16}
\begin{longtable}{@{}>{\raggedright\arraybackslash}p{0.34\linewidth}>{\raggedright\arraybackslash}p{0.22\linewidth}>{\raggedright\arraybackslash}p{0.40\linewidth}@{}}
\toprule
Setting and default & Accepted values & Effect \\
\midrule
\endfirsthead
\multicolumn{3}{@{}l}{Prefix: \nolinkurl{bioseq.aligner.minimizer} (continued)} \\
\toprule
Setting and default & Accepted values & Effect \\
\midrule
\endhead
\bottomrule
\endfoot
% GUC: bioseq.aligner.minimizer.algorithm
\nolinkurl{algorithm}\par Default: \nolinkurl{scalar} & Enum\par \nolinkurl{scalar},\allowbreak\ \nolinkurl{simd_minimap2_reference},\allowbreak\ \nolinkurl{two_stack_nthash},\allowbreak\ \nolinkurl{vulkan_minimap_integer},\allowbreak\ \nolinkurl{bounded_nthash_simd} & Selects the minimizer extraction algorithm. \\[4pt]
% GUC: bioseq.aligner.minimizer.kmer_size
\nolinkurl{kmer_size}\par Default: 15 & Integer\par $1\ldots 64$ & Sets the minimizer k-mer size. \\[4pt]
% GUC: bioseq.aligner.minimizer.window_size
\nolinkurl{window_size}\par Default: 10 & Integer\par $1\ldots 1024$ & Sets the minimizer window size. \\[4pt]
\end{longtable}
\endgroup

The window size counts consecutive $k$-mers, not sequence letters. The ordinary
default is $k=15,w=10$. When requested values equal the ordinary
defaults, some request paths adjust them for short sequences. Explicitly
supplying those same values does not disable that adjustment.
\begin{itemize}
  \item \texttt{simd\_minimap2\_reference} uses scalar extraction for alphabets
    without two-bit ranks.
  \item \texttt{two\_stack\_nthash} requires a nucleotide alphabet, supported
    packed representation, $k\leq31$ and $w\leq256$.
  \item \texttt{bounded\_nthash\_simd} currently requires DNA4, $k=15$ and
    $w=10$. Unsupported requests are rejected.
  \item A Vulkan spelling is accepted independently of whether the build and
    device can execute the required GPU path. See the Vulkan group below.
\end{itemize}

\Needspace{16\baselineskip}
\subsection{Seed-frequency filtering}

Prefix: \nolinkurl{bioseq.aligner.seeding}.

\begingroup
\small
\setlength{\tabcolsep}{4pt}
\renewcommand{\arraystretch}{1.16}
\begin{longtable}{@{}>{\raggedright\arraybackslash}p{0.34\linewidth}>{\raggedright\arraybackslash}p{0.22\linewidth}>{\raggedright\arraybackslash}p{0.40\linewidth}@{}}
\toprule
Setting and default & Accepted values & Effect \\
\midrule
\endfirsthead
\multicolumn{3}{@{}l}{Prefix: \nolinkurl{bioseq.aligner.seeding} (continued)} \\
\toprule
Setting and default & Accepted values & Effect \\
\midrule
\endhead
\bottomrule
\endfoot
% GUC: bioseq.aligner.seeding.maximum_occurrence
\nolinkurl{maximum_occurrence}\par Default: 500 & Integer\par $1\ldots 16777216$ & Limits the target occurrences allowed for one minimizer hash. \\[4pt]
% GUC: bioseq.aligner.seeding.intermediate_occurrence_threshold
\nolinkurl{intermediate_occurrence_threshold}\par Default: 500 & Integer\par $1\ldots 16777216$ & Reserved intermediate occurrence threshold. Stored, but currently unused by execution. \\[4pt]
\end{longtable}
\endgroup

A seed occurrence is a position stored for a minimizer hash. Frequent hashes
can generate many candidate anchors. The occurrence limit controls this
filtering. The intermediate threshold is registered but currently has no
identified execution consumer.

\Needspace{16\baselineskip}
\subsection{Chaining}

Prefix: \nolinkurl{bioseq.aligner.chaining}.

\begingroup
\small
\setlength{\tabcolsep}{4pt}
\renewcommand{\arraystretch}{1.16}
\begin{longtable}{@{}>{\raggedright\arraybackslash}p{0.34\linewidth}>{\raggedright\arraybackslash}p{0.22\linewidth}>{\raggedright\arraybackslash}p{0.40\linewidth}@{}}
\toprule
Setting and default & Accepted values & Effect \\
\midrule
\endfirsthead
\multicolumn{3}{@{}l}{Prefix: \nolinkurl{bioseq.aligner.chaining} (continued)} \\
\toprule
Setting and default & Accepted values & Effect \\
\midrule
\endhead
\bottomrule
\endfoot
% GUC: bioseq.aligner.chaining.algorithm
\nolinkurl{algorithm}\par Default: \nolinkurl{scalar_minimap2} & Enum\par \nolinkurl{scalar_minimap2},\allowbreak\ \nolinkurl{simd_minimap2},\allowbreak\ \nolinkurl{range_minimum_query},\allowbreak\ \nolinkurl{vulkan_tiled_dynamic_programming} & Selects the seed-anchor chaining algorithm. \\[4pt]
% GUC: bioseq.aligner.chaining.maximum_gap
\nolinkurl{maximum_gap}\par Default: 5000 & Integer\par $1\ldots 1073741824$ & Requested chaining gap limit. Scalar chaining uses at least the bandwidth. \\[4pt]
% GUC: bioseq.aligner.chaining.maximum_reference_gap
\nolinkurl{maximum_reference_gap}\par Default: -1 & Integer\par $-1\ldots 1073741824$ & Reserved reference-gap override. Stored, but currently unused by execution. Default -1 denotes the ordinary gap setting. \\[4pt]
% GUC: bioseq.aligner.chaining.bandwidth
\nolinkurl{bandwidth}\par Default: 500 & Integer\par $1\ldots 1048576$ & Bounds the difference between query and target gaps during chaining. Also used by alignment paths. \\[4pt]
% GUC: bioseq.aligner.chaining.long_join_bandwidth
\nolinkurl{long_join_bandwidth}\par Default: 0 & Integer\par $0\ldots 1048576$ & Reserved long-join bandwidth. Stored, but currently unused by execution. \\[4pt]
% GUC: bioseq.aligner.chaining.maximum_skipped_anchors
\nolinkurl{maximum_skipped_anchors}\par Default: 25 & Integer\par $0\ldots 1048576$ & Controls when the selected chaining algorithm stops predecessor search early. \\[4pt]
% GUC: bioseq.aligner.chaining.maximum_predecessors
\nolinkurl{maximum_predecessors}\par Default: 5000 & Integer\par $1\ldots 4194304$ & Bounds the backward predecessor window in dynamic-programming chaining. RMQ does not use it. \\[4pt]
% GUC: bioseq.aligner.chaining.minimum_seed_count
\nolinkurl{minimum_seed_count}\par Default: 3 & Integer\par $0\ldots 1048576$ & Sets the minimum seed count accepted for a chain. \\[4pt]
% GUC: bioseq.aligner.chaining.minimum_score
\nolinkurl{minimum_score}\par Default: 40 & Integer\par $0\ldots 1048576$ & Sets the minimum accepted chain score. \\[4pt]
% GUC: bioseq.aligner.chaining.gap_penalty_scale
\nolinkurl{gap_penalty_scale}\par Default: 0.8 & Real\par $0\ldots 100$ & Scales the chaining gap penalty. \\[4pt]
% GUC: bioseq.aligner.chaining.skip_penalty_scale
\nolinkurl{skip_penalty_scale}\par Default: 0.0 & Real\par $0\ldots 100$ & Scales the chaining skip penalty. \\[4pt]
\end{longtable}
\endgroup

These settings control anchor chaining, described in the chaining chapter.
Distance limits count sequence positions. Scores and penalty scales belong
to the selected chaining rule. Scalar chaining can check one additional
candidate outside the bounded predecessor window.

In SQL, \texttt{simd\_minimap2} chooses the bounded batch approximation. It
requires a supported SIMD target and a positive
\texttt{maximum\_skipped\_anchors}, even though \texttt{SET} accepts zero.
Its output need not match scalar chaining exactly. The reference-gap override
and long-join bandwidth are stored in options but do not currently change
execution.

\Needspace{16\baselineskip}
\subsection{Primary and secondary results}

Prefix: \nolinkurl{bioseq.aligner.secondary}.

\begingroup
\small
\setlength{\tabcolsep}{4pt}
\renewcommand{\arraystretch}{1.16}
\begin{longtable}{@{}>{\raggedright\arraybackslash}p{0.34\linewidth}>{\raggedright\arraybackslash}p{0.22\linewidth}>{\raggedright\arraybackslash}p{0.40\linewidth}@{}}
\toprule
Setting and default & Accepted values & Effect \\
\midrule
\endfirsthead
\multicolumn{3}{@{}l}{Prefix: \nolinkurl{bioseq.aligner.secondary} (continued)} \\
\toprule
Setting and default & Accepted values & Effect \\
\midrule
\endhead
\bottomrule
\endfoot
% GUC: bioseq.aligner.secondary.overlap_threshold
\nolinkurl{overlap_threshold}\par Default: 0.5 & Real\par $0\ldots 1$ & Controls primary/secondary classification using overlap and uncovered query positions. \\[4pt]
% GUC: bioseq.aligner.secondary.uncovered_length_limit
\nolinkurl{uncovered_length_limit}\par Default: 2147483647 & Integer\par $0\ldots 2147483647$ & Maximum uncovered query length used when classifying secondary alignments. \\[4pt]
% GUC: bioseq.aligner.secondary.score_ratio_threshold
\nolinkurl{score_ratio_threshold}\par Default: 0.8 & Real\par $0\ldots 1$ & Ratio used by the secondary-result filter. Zero bypasses this filter. \\[4pt]
% GUC: bioseq.aligner.secondary.maximum_count
\nolinkurl{maximum_count}\par Default: 5 & Integer\par $0\ldots 1024$ & Maximum secondary results retained across one filtering-stage input list. \\[4pt]
\end{longtable}
\endgroup

A secondary result is an alternative alignment classified against a primary
alignment. The overlap is the number of positions shared by the candidate
and primary query intervals. The shorter and longer spans are the lengths of
those two intervals. Here, \emph{uncovered} means query positions in the candidate span
that the current primary alignments do not cover. Classification compares
\[
  \frac{\text{overlap}}{\text{shorter span}}
  -\frac{\text{uncovered}}{\text{longer span}}
  > \texttt{overlap\_threshold}
\]
and also applies the uncovered-length limit. The overlap threshold is therefore
not simply a percentage of shared query bases.

The secondary filter can retain a result by the score ratio or by a score
difference no larger than $2k$. Its count limit applies across the complete
list passed to that filtering stage, not separately to every primary.
A zero score-ratio threshold bypasses both this filter and its count limit.

\Needspace{16\baselineskip}
\subsection{Base-alignment scoring}

Prefix: \nolinkurl{bioseq.aligner.scoring}.

\begingroup
\small
\setlength{\tabcolsep}{4pt}
\renewcommand{\arraystretch}{1.16}
\begin{longtable}{@{}>{\raggedright\arraybackslash}p{0.34\linewidth}>{\raggedright\arraybackslash}p{0.22\linewidth}>{\raggedright\arraybackslash}p{0.40\linewidth}@{}}
\toprule
Setting and default & Accepted values & Effect \\
\midrule
\endfirsthead
\multicolumn{3}{@{}l}{Prefix: \nolinkurl{bioseq.aligner.scoring} (continued)} \\
\toprule
Setting and default & Accepted values & Effect \\
\midrule
\endhead
\bottomrule
\endfoot
% GUC: bioseq.aligner.scoring.match_score
\nolinkurl{match_score}\par Default: 2 & Integer\par $1\ldots 127$ & Sets the alignment match score. \\[4pt]
% GUC: bioseq.aligner.scoring.mismatch_penalty
\nolinkurl{mismatch_penalty}\par Default: 4 & Integer\par $1\ldots 127$ & Sets the alignment mismatch penalty. \\[4pt]
% GUC: bioseq.aligner.scoring.short_gap_open_penalty
\nolinkurl{short_gap_open_penalty}\par Default: 4 & Integer\par $0\ldots 127$ & Sets the short-gap opening penalty. \\[4pt]
% GUC: bioseq.aligner.scoring.short_gap_extension_penalty
\nolinkurl{short_gap_extension_penalty}\par Default: 2 & Integer\par $0\ldots 127$ & Sets the short-gap extension penalty. \\[4pt]
% GUC: bioseq.aligner.scoring.long_gap_open_penalty
\nolinkurl{long_gap_open_penalty}\par Default: 24 & Integer\par $0\ldots 127$ & Sets the long-gap opening penalty. \\[4pt]
% GUC: bioseq.aligner.scoring.long_gap_extension_penalty
\nolinkurl{long_gap_extension_penalty}\par Default: 1 & Integer\par $0\ldots 127$ & Sets the long-gap extension penalty. \\[4pt]
% GUC: bioseq.aligner.scoring.ambiguous_base_penalty
\nolinkurl{ambiguous_base_penalty}\par Default: 1 & Integer\par $0\ldots 127$ & Sets the penalty for ambiguous bases. \\[4pt]
% GUC: bioseq.aligner.scoring.z_drop_threshold
\nolinkurl{z_drop_threshold}\par Default: 400 & Integer\par $-1\ldots 1048576$ & Base-alignment drop threshold. Interpretation depends on the selected aligner. \\[4pt]
% GUC: bioseq.aligner.scoring.end_bonus
\nolinkurl{end_bonus}\par Default: -1 & Integer\par $-1\ldots 1048576$ & KSW2 sequence-end bonus. -1 disables the bonus. \\[4pt]
\end{longtable}
\endgroup

The match value adds to the alignment score. Penalties reduce it. KSW2 can use
the short- and long-gap pairs. Native BlockAligner uses the short-gap pair only:
both penalties must be positive and their sum must not exceed 128. The long-gap
pair and end bonus do not change native BlockAligner execution.

For BlockAligner extensions, \texttt{z\_drop\_threshold} is used as an X-drop
threshold and must be nonnegative. The registered value $-1$ is not a supported
disabling request for that path. Use the base-alignment chapter to interpret
the drop rules rather than treating the names as interchangeable.

\Needspace{16\baselineskip}
\subsection{Base-alignment implementation and block size}

Prefix: \nolinkurl{bioseq.aligner}.

\begingroup
\small
\setlength{\tabcolsep}{4pt}
\renewcommand{\arraystretch}{1.16}
\begin{longtable}{@{}>{\raggedright\arraybackslash}p{0.34\linewidth}>{\raggedright\arraybackslash}p{0.22\linewidth}>{\raggedright\arraybackslash}p{0.40\linewidth}@{}}
\toprule
Setting and default & Accepted values & Effect \\
\midrule
\endfirsthead
\multicolumn{3}{@{}l}{Prefix: \nolinkurl{bioseq.aligner} (continued)} \\
\toprule
Setting and default & Accepted values & Effect \\
\midrule
\endhead
\bottomrule
\endfoot
% GUC: bioseq.aligner.base_alignment.algorithm
\nolinkurl{base_alignment.algorithm}\par Default: \nolinkurl{block_aligner} & Enum\par \nolinkurl{ksw2},\allowbreak\ \nolinkurl{block_aligner} & Selects the base-level alignment algorithm. \\[4pt]
% GUC: bioseq.aligner.block_alignment.minimum_block_size
\nolinkurl{block_alignment.minimum_block_size}\par Default: 32 & Integer\par $8\ldots 4096$ & Sets the minimum block size used by the block aligner. \\[4pt]
% GUC: bioseq.aligner.block_alignment.maximum_block_size
\nolinkurl{block_alignment.maximum_block_size}\par Default: 256 & Integer\par $8\ldots 16384$ & Sets the maximum block size used by the block aligner. \\[4pt]
\end{longtable}
\endgroup

The base-alignment stage runs when requested results require it. For
BlockAligner, sizes are widened to at least the SIMD lane count. The resulting
sizes must be powers of two, and the minimum must not exceed the maximum.
Some invalid combinations reach an internal abort, so a value accepted by
\texttt{SET} is not assurance of a clean SQL error or a valid alignment run.

\Needspace{16\baselineskip}
\subsection{Saved fragment readers and profiling}

Prefix: \nolinkurl{bioseq.aligner}.

\begingroup
\small
\setlength{\tabcolsep}{4pt}
\renewcommand{\arraystretch}{1.16}
\begin{longtable}{@{}>{\raggedright\arraybackslash}p{0.34\linewidth}>{\raggedright\arraybackslash}p{0.22\linewidth}>{\raggedright\arraybackslash}p{0.40\linewidth}@{}}
\toprule
Setting and default & Accepted values & Effect \\
\midrule
\endfirsthead
\multicolumn{3}{@{}l}{Prefix: \nolinkurl{bioseq.aligner} (continued)} \\
\toprule
Setting and default & Accepted values & Effect \\
\midrule
\endhead
\bottomrule
\endfoot
% GUC: bioseq.aligner.fragment_minimizer.lookup_route
\nolinkurl{fragment_minimizer.lookup_route}\par Default: \nolinkurl{control_block_scan} & Enum\par \nolinkurl{control_block_scan},\allowbreak\ \nolinkurl{direct_compact_hash_lookup} & Selects control-block scanning or compact hash-directory lookup for saved fragments. \\[4pt]
% GUC: bioseq.aligner.grouped_soa_generation.enabled
\nolinkurl{grouped_soa_generation.enabled}\par Default: \texttt{off} & Boolean\par \texttt{off}, \texttt{on} & Builds minimizer data grouped by hash, with each field in a separate array. \\[4pt]
% GUC: bioseq.aligner.profiling.enabled
\nolinkurl{profiling.enabled}\par Default: \texttt{off} & Boolean\par \texttt{off}, \texttt{on} & Enables runtime profiling for the aligner pipeline. \\[4pt]
\end{longtable}
\endgroup

The direct lookup uses a compact hash directory to find the saved sequence
positions for a minimizer hash in a fragment reader representation. It requires that compatible image and
its direct-lookup section. Missing required backing causes failure. This is
not a PostgreSQL hash index.

Grouped structure-of-arrays (SoA) generation changes what a supported index
build produces. It is not the switch that enables reading an already-saved
mapped generation. Controlled grouped minimizer rebuilds require both
asynchronous-persistence switches off. Ordinary per-tuple
\texttt{UPDATE\_ON\_INSERT} minimizer maintenance is rejected in this mode.
Streamed loading has a separate post-flush batched-build path. The persistent
index chapter describes the reader image and its lifetime.

\Needspace{16\baselineskip}
\subsection{Vulkan execution}

Prefix: \nolinkurl{bioseq.aligner.vulkan}.

\begingroup
\small
\setlength{\tabcolsep}{4pt}
\renewcommand{\arraystretch}{1.16}
\begin{longtable}{@{}>{\raggedright\arraybackslash}p{0.34\linewidth}>{\raggedright\arraybackslash}p{0.22\linewidth}>{\raggedright\arraybackslash}p{0.40\linewidth}@{}}
\toprule
Setting and default & Accepted values & Effect \\
\midrule
\endfirsthead
\multicolumn{3}{@{}l}{Prefix: \nolinkurl{bioseq.aligner.vulkan} (continued)} \\
\toprule
Setting and default & Accepted values & Effect \\
\midrule
\endhead
\bottomrule
\endfoot
% GUC: bioseq.aligner.vulkan.enabled
\nolinkurl{enabled}\par Default: \texttt{on} & Boolean\par \texttt{off}, \texttt{on} & Permits Vulkan-backed algorithms when their selected route is supported. \\[4pt]
% GUC: bioseq.aligner.vulkan.seed_lookup.enabled
\nolinkurl{seed_lookup.enabled}\par Default: \texttt{off} & Boolean\par \texttt{off}, \texttt{on} & Enables Vulkan seed-hit lookup. \\[4pt]
% GUC: bioseq.aligner.vulkan.index_cache.enabled
\nolinkurl{index_cache.enabled}\par Default: \texttt{on} & Boolean\par \texttt{off}, \texttt{on} & Allows reuse of GPU-resident minimizer-index uploads after a GPU seed path is selected. \\[4pt]
% GUC: bioseq.aligner.vulkan.fused_frontend.enabled
\nolinkurl{fused_frontend.enabled}\par Default: \texttt{off} & Boolean\par \texttt{off}, \texttt{on} & Enables the resident fused Vulkan frontend. \\[4pt]
% GUC: bioseq.aligner.vulkan.index_build.enabled
\nolinkurl{index_build.enabled}\par Default: \texttt{off} & Boolean\par \texttt{off}, \texttt{on} & Builds the target minimizer index on the GPU. \\[4pt]
% GUC: bioseq.aligner.vulkan.alignment_extension.enabled
\nolinkurl{alignment_extension.enabled}\par Default: \texttt{off} & Boolean\par \texttt{off}, \texttt{on} & Allows the GPU base-alignment extension stage in a supported Vulkan route. \\[4pt]
% GUC: bioseq.aligner.vulkan.ordered_sketch_tile_bp
\nolinkurl{ordered_sketch_tile_bp}\par Default: 2048 & Integer\par $1\ldots 16777216$ & Reserved tile size in base pairs. Stored, but currently unused by execution. \\[4pt]
% GUC: bioseq.aligner.vulkan.radix_sort_digit_bits
\nolinkurl{radix_sort_digit_bits}\par Default: 8 & Integer\par $1\ldots 16$ & Reserved radix digit width in bits. Stored, but currently unused by execution. \\[4pt]
\end{longtable}
\endgroup

Enabling Vulkan only permits its use. It does not select a GPU algorithm.
The complete biological GPU route needs coordinated minimizer, chaining,
seed-lookup, fused-frontend and index-build choices, a compatible build and
device, and supported requested results. Turning on one component does not
create an arbitrary CPU/GPU combination.

GPU extension is required for CIGAR output on that route. Upload caching applies
only after a GPU seed path has been selected. The tile-size and radix-digit
settings are registered and copied into options, but no active execution use
was found in the studied implementation.

\Needspace{16\baselineskip}
\subsection{Query-result cache}

Prefix: \nolinkurl{bioseq_cache.query_cache}.

\begingroup
\small
\setlength{\tabcolsep}{4pt}
\renewcommand{\arraystretch}{1.16}
\begin{longtable}{@{}>{\raggedright\arraybackslash}p{0.34\linewidth}>{\raggedright\arraybackslash}p{0.22\linewidth}>{\raggedright\arraybackslash}p{0.40\linewidth}@{}}
\toprule
Setting and default & Accepted values & Effect \\
\midrule
\endfirsthead
\multicolumn{3}{@{}l}{Prefix: \nolinkurl{bioseq_cache.query_cache} (continued)} \\
\toprule
Setting and default & Accepted values & Effect \\
\midrule
\endhead
\bottomrule
\endfoot
% GUC: bioseq_cache.query_cache.mode
\nolinkurl{mode}\par Default: \nolinkurl{disabled} & Enum\par \nolinkurl{disabled},\allowbreak\ \nolinkurl{exact_query},\allowbreak\ \nolinkurl{reusable_fragments} & Disabled, exact-query reuse with request locking, or reuse of completed query-result fragments. \\[4pt]
% GUC: bioseq_cache.query_cache.completed_result_ttl_ms
\nolinkurl{completed_result_ttl_ms}\par Default: 3600000 & Integer\par $0\ldots 2147483647$ & Maximum age of a completed cached result, measured from completion. Unit: ms. \\[4pt]
\end{longtable}
\endgroup

Any non-disabled mode requires \texttt{bioseq\_cache} in
\texttt{shared\_preload\_libraries} and a PostgreSQL restart. Changing the mode
also invalidates cached plans when a transaction is active. The cache is
experimental and may return incorrect results for cases beyond the tested
coverage. The current BioSeqDB cache adapter supports exact reuse only.
Selecting \texttt{reusable\_fragments} does not add partial-range reuse to that
adapter. Section~\ref{sec:appendix-query-cache} explains the distinction.

TTL means \emph{time to live}. The age starts at completion, and reuse takes the
smaller of the producer's TTL and the current reader's TTL. Zero prevents
ordinary reuse after completion. A waiter already registered against that
running computation can still consume its just-completed result once. The
TTL does not stop a long-running query.

\Needspace{16\baselineskip}
\subsection{Controls present only in test builds}

Prefix: \nolinkurl{bioseq.testing}.

\begingroup
\small
\setlength{\tabcolsep}{4pt}
\renewcommand{\arraystretch}{1.16}
\begin{longtable}{@{}>{\raggedright\arraybackslash}p{0.34\linewidth}>{\raggedright\arraybackslash}p{0.22\linewidth}>{\raggedright\arraybackslash}p{0.40\linewidth}@{}}
\toprule
Setting and default & Accepted values & Effect \\
\midrule
\endfirsthead
\multicolumn{3}{@{}l}{Prefix: \nolinkurl{bioseq.testing} (continued)} \\
\toprule
Setting and default & Accepted values & Effect \\
\midrule
\endhead
\bottomrule
\endfoot
% GUC: bioseq.testing.chunked_scanner.enabled
\nolinkurl{chunked_scanner.enabled}\par Default: \texttt{off}\par \textbf{superuser, hidden} & Boolean\par \texttt{off}, \texttt{on} & Routes sequence input through the synthetic chunked scanner. \\[4pt]
% GUC: bioseq.testing.in_memory_scanner.enabled
\nolinkurl{in_memory_scanner.enabled}\par Default: \texttt{off}\par \textbf{superuser, hidden} & Boolean\par \texttt{off}, \texttt{on} & Materializes sequence input before scanner dispatch. \\[4pt]
% GUC: bioseq.testing.persisted_index.failure_mode
\nolinkurl{persisted_index.failure_mode}\par Default: \nolinkurl{none}\par \textbf{superuser, hidden} & Enum\par \nolinkurl{none},\allowbreak\ \nolinkurl{unsupported_vector_shape} & Injects the selected persisted-index adapter failure. \\[4pt]
\end{longtable}
\endgroup

These three controls exist only in builds that enable test indexes. They are
hidden and require superuser permission or the corresponding granted
\texttt{SET} privilege. They are absent from an ordinary production build and
are not recommended search settings.

